(1) 板间视角因子为1，漫灰体辐射度满足 \(J_1=\varepsilon_1\sigma T_1^4+(1-\varepsilon_1)J_2\)，\(J_2=\varepsilon_2\sigma T_2^4+(1-\varepsilon_2)J_1\)。这是各板发射功率与对面入射反射的自洽方程。

(2) 解线性联立式得 \(q=J_1-J_2=\sigma(T_1^4-T_2^4)/[1/\varepsilon_1+1/\varepsilon_2-1]\)，可视作两个表面反射热阻加板间真空传输热阻。\(\varepsilon_1=\varepsilon_2=1\) 时分母为1，恢复黑体 \(q=\sigma(T_1^4-T_2^4)\)。

(3) 热板每面积每秒失热 \(q\)，冷板吸热 \(q\)，所以合计熵产 \(\dot s_{\rm tot}=q(1/T_2-1/T_1)>0\)，因为热流和逆温差同号。任一发射率趋零，分母增大，净热流趋零而熵产亦趋零。

(4) 利用 \((T_0+\Delta T/2)^4-(T_0-\Delta T/2)^4=4T_0^3\Delta T+O(\Delta T^3)\)，得 \(q=G_{\rm rad}\Delta T+O(\Delta T^3)\)，\(G_{\rm rad}=4\sigma T_0^3/[1/\varepsilon_1+1/\varepsilon_2-1]>0\)。\(\Delta T=0\) 无净热流，即使两个板仍向对方发射非零能流。
